\chapter{Runtime Evidence and Review Reference}
\label{app:runtime-reference}

This appendix expands the evidence model used by the implementation. It is not
intended to describe every database column in the source code. Its purpose is to
show the information that must be available to a reviewer after a run finishes.
For a PFE report, this matters because the project is not only an extraction
script; it is a system for producing reviewable technical evidence.

\section{Run Identity}

The run identity records the target and the final state of the investigation.
It gives the operator one stable object to open, compare, rerun, and cite during
review.

\begin{longtable}{>{\raggedright\arraybackslash\bfseries}p{3.7cm}
                  >{\raggedright\arraybackslash}p{4.3cm}
                  >{\raggedright\arraybackslash}p{5.7cm}}
\caption{Run identity fields used for review.}
\label{tab:runtime-run-identity}\\
\toprule
\rowcolor{accent!8}
Field family & Examples & Review purpose \\
\midrule
\endfirsthead
\toprule
\rowcolor{accent!8}
Field family & Examples & Review purpose \\
\midrule
\endhead
\bottomrule
\endfoot
Run key & Run identifier, creation time, update time. & Lets a reviewer refer
to one investigation without ambiguity. \\
Target & Seed URL, normalized URL, submitted batch row. & Preserves what was
actually tested. \\
Status & Running, completed, failed, partial, skipped. & Shows whether the
workflow reached a final result or stopped early. \\
Route summary & Classification result, selected specialist path, final agent. &
Explains why the system chose a landing, hosting, or embedded path. \\
Evidence counts & Number of streams, screenshots, provider rows, emails, tool
calls. & Gives a quick health check before opening the detailed trace. \\
Cost totals & Model cost, token totals, cached-token totals. & Supports
operational comparison between model and prompt versions. \\
\end{longtable}

\section{Agent Execution Records}

Agent records describe what each specialist was asked to do and what it
returned. They prevent a common debugging problem: seeing a final failure
without knowing which agent failed.

\begin{longtable}{>{\raggedright\arraybackslash\bfseries}p{3.7cm}
                  >{\raggedright\arraybackslash}p{4.3cm}
                  >{\raggedright\arraybackslash}p{5.7cm}}
\caption{Agent execution records and why they are separated.}
\label{tab:runtime-agent-records}\\
\toprule
\rowcolor{accent!8}
Record & Stored information & Why it matters \\
\midrule
\endfirsthead
\toprule
\rowcolor{accent!8}
Record & Stored information & Why it matters \\
\midrule
\endhead
\bottomrule
\endfoot
Classification run & Page type, confidence, route reason, signals used. &
Shows whether a later failure started with a routing mistake. \\
Landing run & Candidate links, match metadata, checked sections, duplicate
handling. & Explains whether the discovery stage covered the visible page. \\
Hosting run & Servers attempted, click strategy, popup handling, player state,
handoff URLs. & Shows whether the agent interacted with the player correctly. \\
Embedded run & Frame URL, player type, play action, network harvest result,
failure reason. & Preserves the final player-context investigation. \\
Provider run & Checked hosts, resolved IPs, provider names, country, lookup
status. & Adds infrastructure context to stream evidence. \\
Email run & Evidence grouped, provider target, draft status, review state. &
Keeps communication output separate from technical extraction. \\
\end{longtable}

The separation is useful during regression testing. If a website stops working,
the engineer can compare the last successful run and the new failed run at the
agent level instead of reading a single combined text output.

\section{Runtime Events}

Runtime events form the timeline of a run. They are intentionally small and
timestamped. Each event should answer one question: what just happened?

\begin{table}[!htbp]
\centering
\begin{tabularx}{\textwidth}{Q{3.7cm}Y}
\toprule
\rowcolor{accent!8}
Event type & Example review value \\
\midrule
Run started & Target URL, selected model, browser profile, batch identifier. \\
Agent started & Agent name, task summary, handoff source. \\
Tool requested & Tool name, input summary, current page URL. \\
Tool completed & Status, duration, output summary, error if any. \\
Screenshot captured & Screenshot path, page URL, reason for capture. \\
Stream found & Stream URL hash or host, source tool, frame or page context. \\
Provider resolved & Host, IP, provider name, country, lookup source. \\
Failure recorded & Failure reason, responsible stage, visible blocker. \\
Run completed & Final status, counts, cost, acceptance rule result. \\
\bottomrule
\end{tabularx}
\caption{Runtime event types used to reconstruct a run.}
\label{tab:runtime-event-types}
\end{table}

The final output can be wrong or incomplete; the timeline helps explain why. If
the stream list is empty, events can show whether the agent never clicked play,
whether the click opened a popup, whether the browser lost focus, or whether the
network harvest ran too early.

\section{Tool Call Review}

Tool calls are reviewed differently from model calls. A model response explains
an intended decision, while a tool call records an attempted browser or provider
operation. For this project, the order of tool calls is often more important
than their count.

\begin{longtable}{>{\raggedright\arraybackslash\bfseries}p{3.5cm}
                  >{\raggedright\arraybackslash}p{4.4cm}
                  >{\raggedright\arraybackslash}p{5.8cm}}
\caption{Tool call review checklist.}
\label{tab:runtime-tool-checklist}\\
\toprule
\rowcolor{accent!8}
Tool family & What to check & Common failure it detects \\
\midrule
\endfirsthead
\toprule
\rowcolor{accent!8}
Tool family & What to check & Common failure it detects \\
\midrule
\endhead
\bottomrule
\endfoot
Inspect & Was the right page profile used: landing, hosting, or embedded? &
The agent used generic page text and missed player controls or match cards. \\
Interact & Which selector, text, coordinate, or JavaScript click was used? &
The click opened an ad page, missed a hidden tab, or used navigation for a
same-page control. \\
Navigate & Did the agent intentionally change URL or accidentally reload state?
& The browser lost the active player or reset a server selection. \\
Screenshot & Why was the screenshot captured? & The screenshot became a waiting
loop instead of an evidence checkpoint. \\
Harvest & Was the player active before harvesting? & The tool returned no media
because it ran before network requests were created. \\
Media state & Did the player report play, pause, buffering, or blocked state? &
The agent assumed playback from page appearance alone. \\
Provider lookup & Which host or stream URL was enriched? & Provider analysis
ran on a landing URL instead of the actual stream host. \\
\end{longtable}

This checklist was used informally during debugging and later shaped the
evaluation tables in Chapter~6. It also explains why tool success rate is not
enough. A tool can return successfully while still being called at the wrong
time.

\section{Model Usage and Cost Review}

Model usage fields make the evaluation reproducible. If two prompt versions have
similar website success, the cheaper and more stable version is usually the
better engineering choice.

\begin{table}[!htbp]
\centering
\begin{tabularx}{\textwidth}{Q{4.0cm}Y}
\toprule
\rowcolor{accent!8}
Cost field & Reason for keeping it \\
\midrule
Model provider and model name & Allows the same run to be compared across model
families and rate-limit policies. \\
Prompt version & Links behavior changes to a specific instruction set. \\
Input tokens & Measures full prompt, state, and tool-context load. \\
Cached input tokens & Shows whether repeated prompt/tool context was reused. \\
Output tokens & Measures answer length and reasoning cost. \\
Cache write tokens & Captures the cost of preparing reusable context when
supported by the provider. \\
Estimated input cost & Separates normal input cost from output and cache cost. \\
Estimated output cost & Explains why longer reasoning can increase run price. \\
Total cost & Supports cost per run and cost per working website. \\
\bottomrule
\end{tabularx}
\caption{Model usage and cost fields kept for evaluation.}
\label{tab:runtime-cost-fields}
\end{table}

The most useful metric during development was not total cost alone. It was cost
per working website. A cheaper model that fails many pages can be more expensive
operationally than a slightly stronger model that completes the evidence path
with fewer retries.

\section{Screenshot Review Checklist}

Screenshots are part of the evidence bundle. They are not only debugging
images. A useful screenshot should explain one state of the investigation:
classification, candidate discovery, player activation, blocker, stream
confirmation, or final failure.

\begin{longtable}{>{\raggedright\arraybackslash\bfseries}p{3.6cm}
                  >{\raggedright\arraybackslash}p{4.6cm}
                  >{\raggedright\arraybackslash}p{5.5cm}}
\caption{Screenshot review checklist.}
\label{tab:screenshot-review-checklist}\\
\toprule
\rowcolor{accent!8}
Screenshot moment & What should be visible & Why it is useful \\
\midrule
\endfirsthead
\toprule
\rowcolor{accent!8}
Screenshot moment & What should be visible & Why it is useful \\
\midrule
\endhead
\bottomrule
\endfoot
Classification & Page title area, match list, player shell, or iframe context. &
Supports the selected page type. \\
Landing extraction & Event cards, links, tabs, dates, or channel blocks. &
Shows that candidates came from visible page content. \\
Hosting activation & Server buttons, play overlay, player container, popup state.
& Shows whether the agent interacted with the right controls. \\
Embedded inspection & Player-only page, frame URL context, playback area. &
Supports fallback extraction from iframe/player context. \\
Failure state & Error message, blocked page, empty player, unavailable server. &
Makes failed runs reviewable rather than vague. \\
Final evidence & Player state or network-supported stream context. & Gives the
operator a visual anchor for the final claim. \\
\end{longtable}

The screenshot policy avoids two extremes. Capturing no screenshots makes a run
difficult to review. Capturing screenshots after every small wait creates noise
and cost. The final design captures screenshots at meaningful state changes.

\section{Provider Review Fields}

Provider enrichment links stream evidence to infrastructure context. The lookup
result should be stored even when it is incomplete, because partial information
can still help a reviewer.

\begin{table}[!htbp]
\centering
\begin{tabularx}{\textwidth}{Q{3.7cm}Y}
\toprule
\rowcolor{accent!8}
Provider field & Review meaning \\
\midrule
Original URL & The stream, embedded, or host URL used as lookup input. \\
Hostname & Domain extracted from the URL. \\
Resolved IP & Network address found through DNS resolution. \\
Provider name & Hosting provider, CDN, network owner, or organization. \\
Country & Country associated with the IP or registration data. \\
Lookup source & IP database, RDAP, WHOIS, or other resolver path. \\
Abuse contact & Email or contact field when available. \\
Lookup status & Resolved, provider found, contact missing, failed, or skipped. \\
\bottomrule
\end{tabularx}
\caption{Provider-enrichment fields used during evidence review.}
\label{tab:provider-review-fields}
\end{table}

Separating these fields avoids misleading results. A provider lookup can resolve
an IP and provider but still fail to find an abuse email. That is not the same
as a total lookup failure.

\section{Operator Review Flow}

The operator review flow is the human side of the system. It turns the stored
records into a practical decision: accept the evidence, rerun the website, mark
the result as blocked, or prepare a draft for later enforcement review.

\begin{longtable}{>{\raggedright\arraybackslash\bfseries}p{3.5cm}
                  >{\raggedright\arraybackslash}p{4.8cm}
                  >{\raggedright\arraybackslash}p{5.5cm}}
\caption{Operator review flow after a run finishes.}
\label{tab:operator-review-flow}\\
\toprule
\rowcolor{accent!8}
Review step & Operator checks & Possible decision \\
\midrule
\endfirsthead
\toprule
\rowcolor{accent!8}
Review step & Operator checks & Possible decision \\
\midrule
\endhead
\bottomrule
\endfoot
Open run summary & Final status, page route, evidence counts, cost. & Continue
review, rerun, or reject obvious failure. \\
Inspect trace & Agent transitions, tool order, errors, timestamps. & Identify
which stage caused the result. \\
Inspect screenshots & Page type, player state, blocker, evidence moment. &
Accept visual support or request rerun. \\
Inspect streams & Host, path, manifest type, timestamp, source frame. & Keep
stream evidence or mark as weak/expired. \\
Inspect providers & Provider name, country, IP, contact availability. & Prepare
provider context or flag missing enrichment. \\
Inspect email draft & Evidence grouping, provider target, wording, attachments. &
Send to human approval or discard. \\
Record decision & Accepted, partial, blocked, false positive, needs rerun. &
Feed future regression and prompt/tool improvement. \\
\end{longtable}

This flow is deliberately conservative. The system prepares evidence and review
context, but final operational decisions remain human-controlled.

\section{Acceptance Checklist}

The following checklist summarizes the minimum evidence needed before a run is
treated as useful for review.

\begin{table}[!htbp]
\centering
\begin{tabularx}{\textwidth}{Q{4.2cm}Y}
\toprule
\rowcolor{accent!8}
Checklist item & Required condition \\
\midrule
Correct page role & The page type is explained by visible content and browser
signals. \\
Controlled routing & The next agent receives a focused handoff with a clear
evidence objective. \\
Meaningful interaction & Player clicks, tab changes, and server attempts are
visible in the trace. \\
Evidence capture & At least one stream, embedded URL, manifest, provider row,
or defensible failure reason is recorded. \\
Screenshot support & Screenshots show the page state behind the final claim or
the blocker. \\
Provider context & Stream or host evidence is enriched when possible, with
lookup failures kept explicit. \\
Cost visibility & Model usage and tool counts are stored for evaluation. \\
Review decision & A human can accept, reject, rerun, or mark the result partial
without reconstructing the run from scratch. \\
\bottomrule
\end{tabularx}
\caption{Minimum acceptance checklist for a reviewable investigation run.}
\label{tab:runtime-acceptance-checklist}
\end{table}

This checklist connects the implementation to the PFE objective. The project is
successful when it makes the evidence path understandable and repeatable, not
only when it prints a stream URL.
